\chapter{Special registers}\label{ch:special-registers}

Lane index, simdgroup index, threadgroup position and grid position come from special registers, read with
\op{14059} into a 32-bit register or with \op{14060} into one half.
Every kernel reads at least one; the tensor path reads the lane index to generate fragment addresses.

Three separate numberings identify a special register:

\begin{enumerate}
\def\labelenumi{\arabic{enumi}.}
\item the \textbf{hardware number} in the read instruction's source field, a lookup that Apple's compiler tables do not contain and
  whose field and values depend on the encoded length (\cref{sec:special-registers});
\item the \textbf{register id} Apple's decoder prints (an LLVM register id with a proposed \code{SR\_*} name), numbered
  differently from the hardware field (\cref{sec:special-registers});
\item the \textbf{slot-29 entry} in the kernel's metadata, which declares the registers a kernel reads in a third,
  non-monotone numbering (\cref{sec:slot-29-metadata-numbering}).
\end{enumerate}

Metal builtins do not map one-to-one onto registers: several are computed from a register the kernel reads, so "the
kernel for builtin B reads register R" does not mean R is B (\cref{sec:special-registers}). Whether a read must be waited on is
\cref{sec:read-latency-whether}.

"SR130" means the special register whose 4-byte read form carries byte 0x82 (130 decimal) in its
source field; MM writes registers this way. "Number" is that byte with bit 7 removed (0x82 \& 0x7f = 2). A "builtin" is a
Metal attribute such as \code{[[thread\_index\_in\_simdgroup]]}.

\section{Register numbers and builtins}\label{sec:special-registers}

The read instruction's source field holds a hardware number. The operand is declared as a special-register
class, but its value fits no base-plus-index rule in any of the project's encoding domains. It was recovered as a lookup by
requiring each field bit to be constant across all instances that print the same register and to differ across registers,
then checking the result is a bijection (decoder, \code{isa/g17-special-register-numbers.txt}). The numbering is the
hardware's own: threads-per-threadgroup x, y, z are consecutive numbers 24, 25, 26, while their LLVM
ids are 55, 57, 59. Field position and numbering depend on the form:

\begin{itemize}
\item 4-byte form (1,116 corpus instances): byte 1 bits 0–6; byte 1 bit 7 is set in every witness and is not part of the
  number;
\item 8-byte form (154 instances): byte 1 bits 0–1 and byte 3 bits 3–6, with a different numbering (lane index is 30, not 2).
  The 8-byte form is \op{14060} (Apple's compiler emits it at /8, MM~\mm{25.73}); Apple's decoder
  reports its byte 1 as the register, so an 8-byte read of the simdgroup index prints 131, not 133 (\code{agxforge/g17/mdgen.py}).
  Key an 8-byte read by the decoder's register name, not by byte 1.
\end{itemize}

Only registers on the observed list may be emitted; nothing predicts the number of a register the corpus never read
(\emph{this compiler}, \code{agxforge/g17/cc.py}).

\Cref{tab:special-registers} is the reference list of special registers. Builtin is the Metal attribute whose one-builtin compile reads the register (compile-witnessed,
\code{isa/g17-special-registers.json}); the \code{SR\_*} names are proposed from Apple's library vocabulary and are not themselves
witnessed. "Corpus reads" counts decoded reads over 679 objects and the read form used, 32-bit or 16-bit (decoder,
\code{isa/g17-special-registers.toml}). Apple's corpus defines 200 special registers (66 named \code{SR\_*}) and reads 26 of them,
counting flags.

\widetable
\begin{xltabular}{\linewidth}{@{}>{\hsize=0.84\hsize}X>{\hsize=0.62\hsize}XX>{\hsize=1.40\hsize}X>{\hsize=0.74\hsize}X>{\hsize=1.45\hsize}X>{\hsize=0.90\hsize}X@{}}
\caption{Special registers by 4-byte read byte, with their numbers, decoder id, builtin, corpus reads and standing.}\label{tab:special-registers}\\
\toprule
4-byte byte (SR) & Number (4-byte / 8-byte) & Register id, proposed name & Builtin and meaning & Corpus reads & Standing & Source \\
\midrule
\endfirsthead
\tablecontinued{7}
\toprule
4-byte byte (SR) & Number (4-byte / 8-byte) & Register id, proposed name & Builtin and meaning & Corpus reads & Standing & Source \\
\midrule
\endhead
0x82 (130) & 2 / 30 & 45, \code{SR\_SIMD\_ELEM} & \code{thread\_index\_in\_simdgroup}: lane index 0–31 (also what
\code{thread\_index\_in\_quadgroup} compiles to read) & 183, 16-bit & measured: drives tensor address generation in every executed tensor
kernel & MM~\mm{0.8} \\
0x85 (133) & 5 / 31 & 46, \code{SR\_SIMD\_GRP} & \code{simdgroup\_index\_in\_threadgroup} & 4, 16-bit & measured: tile ownership in executed multi-simdgroup kernels; this
compiler masks it \& 3 for up to four simdgroups & MM~\mm{0.8}, MM~\mm{25.93}, \code{agxforge/g17/cc.py} \\
0x94 (148) & 20 / 23 & 37, \code{SR\_PVQUAD} & not witnessed as a builtin & 0 & decoder only & numbers file \\
0x95 (149) & 21 / 24 & 38, \code{SR\_PVSIMD} & lanes before this one; equals the lane index with every lane active & 15, 16-bit & measured with every lane active only & MM~\mm{25.145.1} \\
0x96 (150) & 22 / - & 65, \code{SR\_TVQUAD} & not witnessed & 0 & decoder only & numbers file \\
0x97 (151) & 23 / 41 & 66, \code{SR\_TVSIMD} & active lanes in all; 32 with every lane active & 15, 16-bit & measured with every lane active only & MM~\mm{25.145.1} \\
0x98–0x9a (152–154) & 24–26 / 35–37 & 55, 57, 59, \code{SR\_TG\_X/Y/Z\_SIZE} & \code{threads\_per\_threadgroup} .x .y .z & x 3, y 4, z 3, 16-bit & x: whole-program (sources reading 152 match Apple's GPU
output); y, z compile-witnessed & MM~\mm{25.149.1}, \code{agxforge/g17/mdgen.py} \\
0x9c–0x9e (156–158) & – & 54, 56, 58, \code{SR\_TG\_X/Y/Z} & \code{threadgroup\_position\_in\_grid} .x .y .z & x 543, y 3, z 4, 32-bit & x measured (every tensor body reads it); y, z values compile-witnessed & MM~\mm{0.8}, MM~\mm{25.145.1} \\
0xa0–0xa2 (160–162) & – & 61, 62, 63, \code{SR\_TP\_IN\_GRID\_X/Y/Z} & \code{thread\_position\_in\_grid} .x .y .z & x 73, y 32, z 34, 32-bit & x measured; y, z values compile-witnessed & MM~\mm{0.8}, MM~\mm{25.145.1} \\
0xa4–0xa6 (164–166) & 36, 37, - / 16–18 & 26, 27, 28, \code{SR\_LOCAL\_X/Y/Z} & \code{thread\_position\_in\_threadgroup} .x .y .z & x 1, z 1, 16-bit & x measured; y, z values compile-witnessed & MM~\mm{0.8}, \code{agxforge/g17/cc.py} \\
0xa7 (167) & 39 / 15 & 25, \code{SR\_LIN\_ID} & \code{thread\_index\_in\_threadgroup} & 13, 16-bit & compile-witnessed & json, numbers file \\
0xa8, 0xaa (168, 170) & 40, 42 / 32–34 & 51, 52, 53, \code{SR\_TG\_DISP\_X/Y/Z\_SIZE} & 51 is \code{threadgroups\_per\_grid.x}; y and z have no witnessed
builtin & 3 each, 16-bit & compile-witnessed (x only) & json, \code{tools/g17emu.py} \\
0xac, 0xad (172, 173) & 44, 45 / - & 23, 24, \code{SR\_LIBXDIM/YDIM} & unknown & – & decoder only & numbers file \\
0xea (234) & 106 / - & 71, \code{SR\_VRID\_SET} & unknown & – & decoder only & numbers file \\
\bottomrule
\end{xltabular}

"Values compile-witnessed" means the register number and builtin are witnessed by compilation, but no dispatch with more
than one threadgroup or thread along that axis has run on hardware; on an axis of extent 1 every reading agrees (MM~\mm{25.145.1}).

No register holds \code{quadgroup\_index\_in\_threadgroup} or \code{simdgroups\_per\_threadgroup}; both are
computed. The first reads \code{SR\_SIMD\_ELEM} and \code{SR\_SIMD\_GRP} (slot 29 carries {[}52, 53, 80{]} with its store's grid position,
\code{agxforge/g17/mdgen.py}), and the second expands to the whole threads-per-threadgroup vector, because it is computed from it
(\code{isa/g17-class-bytes.txt}, \code{tools/g17synth.py}). So a one-builtin compile lists the registers its builtin is
computed from, and that list alone does not identify any of them. \emph{Correction:} the early \code{isa/g17-scalar-isa.toml} names bytes 0x82
and 0x98 after these two builtins; they are the lane index, measured by every executed tensor kernel that generates
fragment addresses from it, and threads-per-threadgroup x (MM~\mm{0.8}, \code{isa/g17-special-register-numbers.txt},
\code{agxforge/g17/cc.py}; \cref{app:a}).

\textbf{Not known.} What \code{SR\_PVSIMD} holds when lanes are inactive. Measured with every lane active, \code{SR\_PVSIMD} reads "lanes before
this one" (the lane index) and \code{SR\_TVSIMD} reads 32; a vote kernel built on them matches Apple's \code{simd\_prefix\_exclusive\_sum(1)}
(MM~\mm{25.145.1}). Separately, Apple's divergent \code{simd\_broadcast\_first} reads register 38 (\code{SR\_PVSIMD}) with the 16-bit read as
the first of four instructions, and MM~\mm{25.144.6} describes it there as "the active-lane mask"; that role is inferred from the decoded
sequence, not measured. With any lane inactive, a count of active lanes before this one and a mask are different values;
neither reading may be assumed under partial activity, only the full-mask values. This compiler's metadata
author refuses 149 and 151 in any case, because they have no measured slot-29 entry (MM~\mm{25.115.1}).

The 16-bit read writes only the half it names. With y = 7t +
0xABCD0000 live, the read wrote y's low half with the lane index and y then read 0xABCD0000 \textbar{} lane, three runs identical: the
read does not zero-extend (measured, MM~\mm{25.145.4}, \code{isa/g17-execution-halfwrite-results.json}). The probe was
run for a delivered 2,048-row attention kernel that wrote a register's low half with the 16-bit read, then shifted the
whole register; it could be bit-exact on hardware only if the stale high half read as zero, either through zero-extension
or because a released half reads zero. With zero-extension ruled out, that kernel depended on its released
high half reading 0, which it did in every run but which is not a rule (\cref{sec:release-does}). The
compiler now shifts only the low half (\op{17016}) (MM~\mm{25.145}, MM~\mm{25.145.1}).

Two further constraints apply when emitting reads:

\begin{itemize}
\item The 4-byte read's destination field is four bits (r0-r15), so read destinations compete for the narrow register pool
  (decoder and this compiler, MM~\mm{25.110}; \cref{ch:register-file}).
\item Reading threadgroup position, thread position in threadgroup, \code{SR\_PVSIMD}, \code{SR\_TVSIMD}, \code{SR\_SIMD\_ELEM} or \code{SR\_SIMD\_GRP} clears the
  per-kernel "every thread independent" metadata flag (Apple's compiler emits, 47 one-variable probes, \code{agxforge/g17/cc.py});
  \cref{ch:kernel-metadata-image}.
\end{itemize}

\section{Slot-29 metadata numbering}\label{sec:slot-29-metadata-numbering}

The kernel's metadata (per-kernel slot 29) lists the special registers the kernel reads, but in its own numbering. It is a
separate namespace from the hardware numbers above (MM~\mm{0.8}); the metadata format is \cref{ch:kernel-metadata-image}'s. \Cref{tab:slot-29-entries} gives the known entries.

\begin{xltabular}{\linewidth}{@{}Xl>{\hsize=1.26\hsize}X>{\hsize=0.69\hsize}X@{}}
\caption{Slot-29 metadata entries for the special registers.}\label{tab:slot-29-entries}\\
\toprule
Register (4-byte byte) & Slot-29 entry & Standing & Source \\
\midrule
\endfirsthead
\tablecontinued{4}
\toprule
Register (4-byte byte) & Slot-29 entry & Standing & Source \\
\midrule
\endhead
\code{SR\_SIMD\_ELEM} (130) & 52 & measured for the active tensor form: 1,039 exact (130, 156) objects all
carry {[}0, 52{]} & MM~\mm{0.8}, \code{agxforge/g17/mdgen.py} \\*
\code{SR\_SIMD\_GRP} (133) & 53 & compiled alone it carries {[}53{]}; with grid position {[}53, 80{]}; the
multi-simdgroup tensor class (52, 53) & MM~\mm{25.93}, \code{agxforge/g17/mdgen.py} \\*
threads\_per\_threadgroup.x (152) & 4 & one object reading 152 alone carries {[}4{]}; five sources reading it
with others agree and match Apple on hardware & MM~\mm{25.149.1}, \code{agxforge/g17/mdgen.py} \\*
threadgroup\_position\_in\_grid .x .y .z (156–158) & 0, 1, 2 & measured & MM~\mm{0.8} \\*
thread\_position\_in\_grid .x .y .z (160–162) & 80, 81, 82 & measured & MM~\mm{0.8} \\*
thread\_position\_in\_threadgroup .x .y .z (164–166) & 48, 49, 50 & measured & MM~\mm{0.8} \\*
\bottomrule
\end{xltabular}

\begin{itemize}
\item The bases are a lookup. 156 → 0, 160 → 80 and 164 → 48 are not monotone in the register number, so \code{entry = register - 80}, which fits the middle row, is refuted by both others. Only the axis increment within a base is arithmetic.
\item The SR130 entry depends on the read form. Decoder forms whose first printed register byte is also 130 but whose full
  identity (byte 1, width, byte 3) differs map differently, so no universal 130 → 52 rule exists (MM~\mm{0.8}).
\item For 43 sources the slot-29 law predicted the vector in Apple's own compile exactly in
  37, and the other 6 named 152 or a five-register set, which the later entries cover (MM~\mm{25.149.1}).
\item \emph{Correction:} MM~\mm{0.8}'s table lists \code{SR\_SIMD\_GRP} as "not generalized", with no slot-29 entry; the entry is 53 (MM~\mm{25.93},
  \code{agxforge/g17/mdgen.py}; \cref{app:a}).
\end{itemize}

\section{Read latency and waits}\label{sec:read-latency-whether}

A special-register read publishes a scoreboard slot (slot semantics: \cref{sec:scoreboard}; field positions: 
\cref{sec:scoreboard-publish-wait}), and Apple's compiler waits on it, but a read that publishes on slot 0 is readable by the next instruction with no wait (measured, MM~\mm{25.117}, MM~\mm{25.121}). Apple's
wait is a scheduling choice, and its measured cost is within run-to-run spread.

An early ledger entry recorded that a store reading a special-register value
directly stored 0, and attributed it to the store taking a value that had not yet arrived. This compiler then waited before such stores, and
MM~\mm{6}, MM~\mm{0.8} and MM~\mm{25.90} list the read among late values for that reason. A latency probe later tested that explanation against others (\code{tools/g17srlatency.py}, every arm at one threadgroup and at 8,192 groups, three rounds, predictions committed;
\code{isa/g17-execution-sr-latency-results.json}, 90 of 90 rows reproduced); \cref{tab:sr-latency-arms} gives the arms.

\begin{xltabular}{\linewidth}{@{}>{\hsize=0.58\hsize}X>{\hsize=1.42\hsize}Xl@{}}
\caption{Arms of the special-register latency probe and what each store wrote.}\label{tab:sr-latency-arms}\\
\toprule
Arm & Store & Stored \\
\midrule
\endfirsthead
\tablecontinued{3}
\toprule
Arm & Store & Stored \\
\midrule
\endhead
d0 to d4 & \code{C[t] = t} straight off the read, 0 to 4 independent instructions
between & 0 \\*
d0\_wait & d0 with the store waiting on the read's slot & 0 \\*
sr5 & read writes R5; store value R5, index R4 (distance 0) & t \\*
same\_release & \code{O[u] = u} with u from an ALU, value and index both R16 (no special
register at all) & 0 \\*
same\_op6\_cleared & that store with only the index's lifetime 16
→ 0 & u \\*
\bottomrule
\end{xltabular}

Neither distance nor a wait fixed the failing arms, and distance 0 worked once the value and the index were different
registers. Every failing arm used one register as both the store's value and its index, with the index slot releasing it:
operand 6 = 16 on \op{17229} clears the register before the value is read. The defect was this aliased release
(\cref{ch:register-file}), which breaks any value stored at its own index, whether or not it came from a special register.

Apple's compiler waits uniformly (measured, MM~\mm{25.121}), after all 27 special registers it reads and
at every distance, including 218 first consumers four or more instructions later; 3,496 of 3,497 first consumers at distance 1
wait, and 8,808 of 8,825 first consumers over 6,594 programs (MM~\mm{25.144.8}). A chain of 4,000 reads with
every consumer waiting ran 1.092~ms against 1.096 unwaited at the full grid and 0.1528 against 0.1521~ms at one simdgroup,
inside run-to-run spread, every lane right. The waited encoding is operand 1 bit 24 (byte 1 bit 2) of \op{10279} and \op{10282}. Delivered programs from this compiler read special registers unwaited 6,604 times and are
bit-exact on hardware (MM~\mm{25.145}).

\leadhead{Safe assumptions.}

\begin{itemize}
\item A slot-0 special-register read can be consumed at distance 0 without a wait. The evidence covers \op{14059} values consumed by stores and adds, and both read forms in delivered programs.
\item Another fill slot cannot be assumed to behave the same. The same 4,000-read program with each read publishing on slot 7
  (operand 1 bits 20–23 = 8) stored nothing, waited or not. The read uses the fill slot it is given, and for this read
  slot 7 does not behave like slot 0 (measured, unexplained, MM~\mm{25.121}).
\item A wait on a slot-0 read is harmless if a scheduler prefers uniformity.
\end{itemize}

\emph{Correction:} "a read\_sr fills a scoreboard slot like a load, so its first consumer must wait on it" (MM~\mm{6}, MM~\mm{0.8},
MM~\mm{25.90}) is superseded for slot-0 reads by MM~\mm{25.117} and MM~\mm{25.121}.

\textbf{Evidence.} MM~\mm{0.8}, MM~\mm{6}, MM~\mm{25.73}, MM~\mm{25.90}, MM~\mm{25.93}, MM~\mm{25.110}, MM~\mm{25.115.1}, MM~\mm{25.117}, MM~\mm{25.121}, MM~\mm{25.144.8},
MM~\mm{25.144.6}, MM~\mm{25.145}, MM~\mm{25.145.1}, MM~\mm{25.145.4}, MM~\mm{25.149.1}; \code{isa/g17-special-register-numbers.txt},
\code{isa/g17-special-registers.toml}, \code{isa/g17-special-registers.json}, \code{isa/g17-scalar-isa.toml}, \code{isa/g17-class-bytes.txt},
\code{agxforge/g17/cc.py}, \code{agxforge/g17/mdgen.py}, \code{tools/g17emu.py}, \code{tools/g17synth.py}.
